% DFFRNT AI Assistant — User Guide
% Overleaf-compatible (pdfLaTeX). Styling matches the DFFRNT report theme:
% Palatino body, garnet (uOttawa-inspired) accent, markdown-style callouts,
% and the themed table style for the configuration annex.
\documentclass[11pt,letterpaper]{article}

%------------------------------------------------------------------ encoding --
\usepackage[T1]{fontenc}
\usepackage[utf8]{inputenc}

% Silence the harmless "Command \showhyphens has changed" notice microtype
% emits under recent LaTeX kernels.
\usepackage{silence}
\WarningFilter{latex}{Command}

%--------------------------------------------------------------------- fonts --
\usepackage{mathpazo}                 % Palatino body + matching math
\usepackage[scaled=0.92]{helvet}      % Helvetica-like sans (headings)
\usepackage{microtype}                % refined justification & spacing
\linespread{1.08}
\renewcommand{\sfdefault}{\rmdefault} % unify on Palatino, matching the report

%-------------------------------------------------------------------- layout --
\usepackage[letterpaper,margin=1in,headheight=22pt,headsep=18pt,footskip=28pt]{geometry}
\usepackage{tocloft}                  % load BEFORE parskip
\usepackage{parskip}
\usepackage{enumitem}
\setlist{itemsep=2pt,topsep=4pt,parsep=0pt}
\usepackage{environ}

%--------------------------------------------------------------------- colour --
\usepackage[table]{xcolor}
\definecolor{accent}{HTML}{6A24B5}    % DFFRNT brand purple
\definecolor{accentdk}{HTML}{4E1A86}  % deeper purple (headings / links)
\definecolor{brandlime}{HTML}{CFEF3C} % chartreuse accent
\definecolor{pagebg}{HTML}{E7E2F0}    % lavender title-page background
\definecolor{icnk}{HTML}{171634}      % dark navy icon tile
\definecolor{blocklt}{HTML}{D9CEEE}   % light lavender text on the purple block
\definecolor{blocklt2}{HTML}{C7B9E6}  % lighter lavender (block small print)
\definecolor{ink}{HTML}{1F2933}       % near-black body ink
\definecolor{inkmute}{HTML}{56616C}   % muted grey for metadata
\definecolor{rulegrey}{HTML}{C7CDD4}  % hairline rules
\definecolor{briefbg}{HTML}{F1ECF8}   % light lavender callout / code background

% Admonition palette (markdown-style callouts)
\definecolor{noteblue}{HTML}{2B6CB0}
\definecolor{notebg}{HTML}{EEF4FB}
\definecolor{tipgreen}{HTML}{2F7A55}
\definecolor{tipbg}{HTML}{EDF6F0}
\definecolor{warnred}{HTML}{B02A37}
\definecolor{warnbg}{HTML}{FBEEEF}

%-------------------------------------------------------------------- tables ---
\usepackage{array}
\usepackage{booktabs}
\usepackage{tabularx}
\usepackage{xltabular}   % page-breakable tabularx (loads longtable + tabularx)
\definecolor{tablehead}{HTML}{6A24B5}   % header fill (= accent)
\definecolor{tableshade}{HTML}{F1ECF8}  % alternating-row shade

% Auto-stretching, full-width column types
\newcolumntype{L}{>{\raggedright\arraybackslash}X}
\newcolumntype{C}{>{\centering\arraybackslash}X}

% Header-cell text style (white bold sans on the garnet bar)
\newcommand{\thd}[1]{{\sffamily\bfseries\color{white}#1}}

% ptable: styled non-breaking table body. Argument = column spec.
\NewEnviron{ptable}[1]{%
  \par\vspace{4pt}\footnotesize
  \setlength{\extrarowheight}{3pt}%
  \arrayrulecolor{rulegrey}%
  \rowcolors{2}{tableshade}{white}%
  \begin{tabularx}{\linewidth}{#1}
  \rowcolor{tablehead}\BODY
  \end{tabularx}%
  \arrayrulecolor{black}\par\vspace{6pt}%
}

% pltable: page-breaking sibling of ptable (header declared with \endhead
% inside the body so it repeats on every page).
\NewEnviron{pltable}[1]{%
  \footnotesize
  \setlength{\extrarowheight}{3pt}%
  \arrayrulecolor{rulegrey}%
  \rowcolors{2}{tableshade}{white}%
  \begin{xltabular}{\linewidth}{#1}
  \BODY
  \end{xltabular}%
  \arrayrulecolor{black}%
}

% Re-ingest warning marker for the annex (index-invalidating parameters)
\newcommand{\warnmark}{{\color{warnred}\sffamily\bfseries\,(!)}}

%------------------------------------------------------------- code listings ---
\usepackage{listings}
\lstset{
  basicstyle=\ttfamily\small,
  keywordstyle=\color{accentdk}\bfseries,
  commentstyle=\color{inkmute}\itshape,
  stringstyle=\color{accent},
  frame=single,
  rulecolor=\color{rulegrey}, framesep=6pt,
  showstringspaces=false, breaklines=true, tabsize=2,
  backgroundcolor=\color{briefbg!40},
  xleftmargin=4pt,
}

%---------------------------------------------------- section / TOC styling ---
\usepackage{titlesec}
\titleformat{\section}
  {\sffamily\LARGE\bfseries\color{accentdk}}
  {\makebox[1.6em][l]{\thesection}}{0em}{}
  [\vspace{3pt}{\color{accent}\titlerule[1pt]}]
\titleformat{\subsection}
  {\sffamily\Large\bfseries\color{ink}}
  {\makebox[2.0em][l]{\thesubsection}}{0em}{}
\titleformat{\subsubsection}
  {\sffamily\large\bfseries\color{ink}}
  {\thesubsubsection}{0.6em}{}
% FAQ questions & bolded Q-style leads: garnet, run-in
\titleformat{\paragraph}[runin]
  {\sffamily\bfseries\color{accentdk}}{}{0em}{}
\titlespacing*{\section}      {0pt}{0pt}{14pt}
\titlespacing*{\subsection}   {0pt}{16pt}{6pt}
\titlespacing*{\subsubsection}{0pt}{12pt}{4pt}
\titlespacing*{\paragraph}    {0pt}{8pt}{0.7em}

\renewcommand{\thesection}{\arabic{section}}
\renewcommand*\contentsname{Table of Contents}
\renewcommand{\cfttoctitlefont}{\sffamily\LARGE\bfseries\color{accentdk}}
\renewcommand{\cftsecfont}{\sffamily\bfseries\color{ink}}
\renewcommand{\cftsecpagefont}{\sffamily\bfseries\color{ink}}
\renewcommand{\cftsubsecfont}{\color{inkmute}}
\renewcommand{\cftsubsecpagefont}{\color{inkmute}}
\setlength{\cftbeforesecskip}{6pt}
\renewcommand{\cftsecleader}{\cftdotfill{\cftdotsep}}

%-------------------------------------------------------- header / footer -----
\usepackage{fancyhdr}
\pagestyle{fancy}
\fancyhf{}
\renewcommand{\headrulewidth}{0.4pt}
\renewcommand{\headrule}{\hbox to\headwidth{\color{rulegrey}\leaders\hrule height \headrulewidth\hfill}}
\renewcommand{\footrulewidth}{0pt}
\renewcommand{\sectionmark}[1]{\markboth{\thesection\quad #1}{}}
\fancyhead[L]{\footnotesize\sffamily\color{inkmute}User Guide}
\fancyhead[R]{\footnotesize\sffamily\color{inkmute}\nouppercase{\leftmark}}
\fancyfoot[R]{\footnotesize\sffamily\color{inkmute}Page \thepage}
\fancyfoot[L]{\footnotesize\sffamily\color{inkmute}DFFRNT Air-Gapped AI Assistant}
\fancypagestyle{plain}{%
  \fancyhf{}%
  \renewcommand{\headrulewidth}{0pt}%
  \fancyfoot[R]{\footnotesize\sffamily\color{inkmute}Page \thepage}%
  \fancyfoot[L]{\footnotesize\sffamily\color{inkmute}DFFRNT Air-Gapped AI Assistant}%
}

%------------------------------------------------------- hyperlinks (late) ----
\usepackage[colorlinks=true,linkcolor=accentdk,urlcolor=accent,
            citecolor=accentdk,linktoc=all]{hyperref}
\usepackage{url}
\urlstyle{sf}

%------------------------------------------------ markdown-style callouts ------
\usepackage[most]{tcolorbox}
\usetikzlibrary{shapes.geometric}   % 4-point "sparkle" star on the title page
\usepackage{graphicx}
\graphicspath{{./}{doc/}}
% Lime-tile icon. Uses the real DFFRNT PNG logo when `dffrnt-icon.png` is present
% beside this file; otherwise falls back to a vector sparkle so the page always
% renders. Drop the PNG in to get the true logo; delete this macro's fallback
% branch (leaving just the circle) to have empty tiles instead.
\newcommand{\dffrntlogo}[1]{%
  \IfFileExists{dffrnt-icon.png}
    {\node at (#1) {\includegraphics[width=15mm]{dffrnt-icon.png}};}
    {\node[icontile] at (#1) {};
     \node[sparkle]  at (#1) {};
     \node[star, star points=4, star point ratio=0.34, fill=brandlime,
           minimum size=2.6mm, inner sep=0] at ([xshift=-3.9mm,yshift=-3.4mm]#1) {};}%
}
\newtcolorbox{gnote}{
  enhanced, breakable,
  colback=notebg, colframe=noteblue,
  boxrule=0pt, leftrule=3.5pt, sharp corners,
  left=12pt, right=12pt, top=8pt, bottom=8pt,
  before skip=10pt, after skip=12pt,
  fontupper=\small\color{ink},
  before upper={{\sffamily\bfseries\color{noteblue}Note}\par\smallskip},
}
\newtcolorbox{gtip}{
  enhanced, breakable,
  colback=tipbg, colframe=tipgreen,
  boxrule=0pt, leftrule=3.5pt, sharp corners,
  left=12pt, right=12pt, top=8pt, bottom=8pt,
  before skip=10pt, after skip=12pt,
  fontupper=\small\color{ink},
  before upper={{\sffamily\bfseries\color{tipgreen}Tip}\par\smallskip},
}
\newtcolorbox{gwarn}{
  enhanced, breakable,
  colback=warnbg, colframe=warnred,
  boxrule=0pt, leftrule=3.5pt, sharp corners,
  left=12pt, right=12pt, top=8pt, bottom=8pt,
  before skip=10pt, after skip=12pt,
  fontupper=\small\color{ink},
  before upper={{\sffamily\bfseries\color{warnred}Warning}\par\smallskip},
}

\setcounter{tocdepth}{2}
\setcounter{secnumdepth}{3}

% Start each major section a little lower.
\let\oldsection\section
\renewcommand{\section}{\vspace{1em}\oldsection}

%==============================================================================
\begin{document}
%==============================================================================

%--------------------------------------------------------------- TITLE PAGE ----
\begin{titlepage}
\thispagestyle{empty}
\begin{tikzpicture}[remember picture, overlay,
    icontile/.style={fill=icnk, rounded corners=2mm, minimum size=15mm, inner sep=0},
    sparkle/.style={star, star points=4, star point ratio=0.36, fill=brandlime,
                    minimum size=8mm, inner sep=0}]
  % page background
  \fill[pagebg] (current page.north west) rectangle (current page.south east);

  % --- decorative circle cluster (top) ---
  \fill[accent]    ([xshift=83mm, yshift=-40mm]current page.north west) circle (17mm);
  \fill[accent]    ([xshift=167mm,yshift=-43mm]current page.north west) circle (17mm);
  \fill[accent]    ([xshift=131mm,yshift=-85mm]current page.north west) circle (15mm);
  \fill[accent]    ([xshift=170mm,yshift=-122mm]current page.north west) circle (14mm);
  \coordinate (limeA) at ([xshift=128mm,yshift=-38mm]current page.north west);
  \coordinate (limeB) at ([xshift=172mm,yshift=-86mm]current page.north west);
  \fill[brandlime] (limeA) circle (17mm);
  \fill[brandlime] (limeB) circle (17mm);
  \dffrntlogo{limeA}
  \dffrntlogo{limeB}

  % --- product title + tagline ---
  \node[anchor=west, text=accentdk,
        font=\fontfamily{phv}\bfseries\fontsize{40}{46}\selectfont]
    at ([xshift=26mm,yshift=-152mm]current page.north west) {DFFRNT};
  \node[anchor=west, text=accentdk,
        font=\fontfamily{phv}\bfseries\fontsize{40}{46}\selectfont]
    at ([xshift=26mm,yshift=-168mm]current page.north west) {AI Assistant};
  \node[anchor=north west, text=accent, font=\fontfamily{phv}\selectfont\large]
    at ([xshift=26mm,yshift=-180mm]current page.north west)
    {\parbox{155mm}{A fully air-gapped, private knowledge and content AI assistant}};

  % --- guide block ---
  \fill[accent, rounded corners=5mm]
    ([xshift=20mm,yshift=-190mm]current page.north west)
    rectangle ([xshift=196mm,yshift=-258mm]current page.north west);
  \node[anchor=west, text=brandlime,
        font=\fontfamily{phv}\bfseries\fontsize{42}{46}\selectfont]
    at ([xshift=32mm,yshift=-212mm]current page.north west) {USER GUIDE};
  \node[anchor=north west, text=blocklt, font=\fontfamily{phv}\selectfont\large]
    at ([xshift=32mm,yshift=-222mm]current page.north west)
    {\parbox{152mm}{Managing the stack and working in the assistant}};
  \node[anchor=north west, text=blocklt2, font=\fontfamily{phv}\selectfont\footnotesize]
    at ([xshift=32mm,yshift=-234mm]current page.north west)
    {\parbox{152mm}{Everyday use of the assistant: the dffrnt-manager control
     panel and CLI, the chat interface, the document library and tagging system,
     routine tasks such as swapping models, and a full config.toml reference.}};
  \node[anchor=west, text=brandlime, font=\fontfamily{phv}\bfseries\selectfont\small]
    at ([xshift=32mm,yshift=-252mm]current page.north west) {\today};
\end{tikzpicture}
\end{titlepage}

%---------------------------------------------------------- FRONT MATTER -------
\pagenumbering{roman}
\setcounter{page}{1}
\pagestyle{fancy}

{\hypersetup{linkcolor=ink}\tableofcontents}
\clearpage

%------------------------------------------------------------- MAIN MATTER ------
\pagenumbering{arabic}
\setcounter{page}{1}

This guide covers everyday use of the DFFRNT Air-Gapped AI Assistant:
managing the stack with \texttt{dffrnt-manager}, working in the assistant
application itself (chat, document library, tags), performing routine tasks
like swapping models, and --- in the annex --- a full reference for every
\texttt{config.toml} parameter.

For getting the software onto a machine in the first place, see the companion
\emph{Installer Guide}.

\section{The dffrnt-manager application}

\texttt{dffrnt-manager} is the single control surface for the whole stack. It
lives in the install folder (e.g.\ \texttt{dffrnt/}) and has two faces: a
\textbf{browser control panel} and an identical \textbf{command-line
interface}. Everything the panel can do, the CLI can do, and vice versa.

\subsection{Launching}

\textbf{Browser panel} --- from the install folder:

\begin{lstlisting}
cd dffrnt
./dffrnt-manager
\end{lstlisting}

This starts a small local web server and opens the panel in your default
browser. The URL looks odd on purpose: it serves on \texttt{127.0.0.1} (this
machine only) and carries a random per-session token, so nothing on the
network --- and no other user on the machine --- can reach the management
functions.

\textbf{Remote / headless machines} --- the panel can't open a browser over
SSH, so print the URL instead and tunnel to it:

\begin{lstlisting}
./dffrnt-manager panel --no-browser     # prints the URL (note the port + token)
# on your own machine:
ssh -L 8901:127.0.0.1:8901 user@target  # then open the printed URL locally
\end{lstlisting}

\textbf{CLI} --- any argument makes it run headless:

\begin{lstlisting}
./dffrnt-manager status
\end{lstlisting}

\subsection{The panel views}

\begin{ptable}{>{\sffamily\bfseries}p{2.2cm} L}
\thd{View} & \thd{What it does}\\
Install & First-run deployment: pick the \texttt{dffrnt-*.tar.gz} bundle
(auto-detected when next to the manager) and a destination, pass the
prerequisite check, unpack, load images, first start. After installation
you'll normally never need this view again.\\
Manage & Start / stop / restart the stack, live service + API health, and
force a re-ingest of every stored document.\\
Models & See the model store (names, sizes, roles), switch the active LLM,
and move models on/off an air-gapped machine as tarballs (import/export).\\
Logs & Follow the containers' logs live (all services or one at a time) ---
the first place to look when something misbehaves.\\
\end{ptable}

Only one operation runs at a time --- if a button reports ``busy'', another
job (perhaps started from a second browser tab) is still in progress. Job
output survives a page reload; you can close the tab and come back
mid-install.

\subsection{Command reference}

All commands run from the install folder.

\begin{pltable}{>{\ttfamily\footnotesize}p{6.4cm} L}
\rowcolor{tablehead}\thd{Command} & \thd{What it does}\\
\endfirsthead
\rowcolor{tablehead}\thd{Command} & \thd{What it does}\\
\endhead
./dffrnt-manager & \normalfont Serve + open the browser panel\\
./dffrnt-manager panel --no-browser & \normalfont Serve the panel, print the
URL only (SSH tunnels); \texttt{--port N} fixes the port\\
./dffrnt-manager install [BUNDLE] [--dest DIR] [--overwrite-config] &
\normalfont Deploy a bundle (auto-detects one next to the binary)\\
./dffrnt-manager start & \normalfont Bring the whole stack up (Qdrant +
Ollama + API)\\
./dffrnt-manager stop & \normalfont Stop all services\\
./dffrnt-manager restart & \normalfont Stop, then start --- required after
editing \texttt{config.toml}\\
./dffrnt-manager status & \normalfont Per-service + API health; exits
non-zero if degraded (script-friendly)\\
./dffrnt-manager logs [all|qdrant|ollama|api] & \normalfont Follow service
logs (Ctrl-C to stop)\\
./dffrnt-manager audit & \normalfont Follow the application audit trail
(\texttt{logs/audit.jsonl})\\
./dffrnt-manager reingest [--dry-run] [--yes] & \normalfont Force re-ingest
of every stored document (previews first)\\
./dffrnt-manager models & \normalfont List the model store with sizes and
roles\\
./dffrnt-manager models use NAME & \normalfont Switch the active LLM and
restart\\
./dffrnt-manager models pull [NAME...] & \normalfont Pull model(s) from the
internet (online machines)\\
./dffrnt-manager models rm NAME & \normalfont Delete a model from the store\\
./dffrnt-manager models export NAME DEST.tar & \normalfont Pack a model into
a tarball (for USB transfer)\\
./dffrnt-manager models import TARBALL & \normalfont Load a model tarball
into the store (checksum-verified)\\
\end{pltable}

\subsection{Manager FAQ}

\paragraph{Do I have to keep the panel open?}
No. The panel is only a remote control --- closing the tab (or the whole
manager) leaves the assistant running. The stack keeps serving until you
\texttt{stop} it.

\paragraph{The panel URL stopped working after I restarted the manager.}
That's expected: the token in the URL is minted per session. Run
\texttt{./dffrnt-manager} again and use the newly printed/opened URL.

\paragraph{Can two people manage at once?}
The panel accepts one mutating operation at a time; a second request gets a
``busy'' response until the first finishes. Status and log views are
read-only and always available.

\paragraph{Where does the manager itself log?}
\texttt{logs/installer.log} in the install folder records every manager
operation; \texttt{./dffrnt-manager logs} follows the \emph{services'} logs
instead.

\paragraph{Does the CLI have anything the panel doesn't?}
\texttt{audit} (following the app audit trail) and \texttt{dev} (a
development mode for source checkouts) are CLI-only. Everything an operator
needs day-to-day is in both.

\subsection{Manager troubleshooting}

\begin{itemize}
  \item \textbf{``busy'' that never clears} --- a long job (model pull,
        re-ingest) is still running; watch it in Logs. If the manager truly
        died mid-job, restart the panel process and check
        \texttt{./dffrnt-manager status}.
  \item \textbf{Panel opens but every action fails} --- usually Docker
        stopped underneath it. Check \texttt{docker ps}; start Docker
        (Desktop or \texttt{systemctl start docker}) and retry.
  \item \textbf{\texttt{status} exits degraded right after \texttt{start}}
        --- services take a moment to pass health checks, and the models take
        longer to warm. Wait a minute and re-run.
  \item For permissions problems, missing bundles, and port conflicts, see
        the Installer Guide's troubleshooting section --- those issues are
        identical post-install.
\end{itemize}

\section{The DFFRNT Assistant application}

Open \url{http://localhost:8000} in a browser (any modern browser on the
machine, or on the network if the port is exposed --- remember the app has no
login). The window has two areas:

\begin{itemize}
  \item a \textbf{sidebar} on the left --- brand, \emph{New chat},
        \emph{Document Library}, a search box, and your recent conversations
        (collapsible via the toggle at the top);
  \item the \textbf{main area} --- either the chat panel or the Document
        Library.
\end{itemize}

The assistant answers \textbf{only from the documents in its library}. If the
documents don't contain the answer, it says so rather than guessing --- that
is deliberate. So the workflow is: upload documents $\rightarrow$ organise
them with tags $\rightarrow$ ask questions.

\subsection{The chat interface}

Type a question into the composer (``Ask anything about your
documents\ldots'') and press \textbf{Enter} (or the send button). The answer
streams in token by token.

\paragraph{Thinking.}
The models are reasoning models: they ``think'' before answering. The
\textbf{Thinking on/off} pill above the composer controls whether that
reasoning is shown as a live block above the answer or folded away behind a
``Thinking\ldots'' label. It's a display choice only --- the model reasons
either way --- and it persists across sessions.

\paragraph{Citations and sources.}
Every claim in an answer carries an inline bracketed citation like
\texttt{[1]} or \texttt{[2][3]}. The numbers map to the source list shown
with the answer; clicking a citation highlights the matching source. This is
your audit trail --- the disclaimer under the composer means what it says:
\emph{verify important information from source documents}.

\paragraph{Message actions.}
Hovering a message reveals its actions:

\begin{itemize}
  \item \textbf{Copy} --- copy the message text to the clipboard.
  \item \textbf{Retry} --- regenerate the assistant's answer to the same
        question (useful after changing scope or uploading more documents).
  \item \textbf{Download} --- save an answer as \textbf{Plain text (.txt)},
        \textbf{Markdown (.md)}, or \textbf{PDF (.pdf)}.
  \item \textbf{Stop} --- while an answer is streaming, stop generation.
\end{itemize}

\paragraph{Refusals.}
``The available documents do not contain enough information to answer
this.''\ is the assistant declining to guess. If you believe the answer
\emph{is} in the library, rephrase with the document's own vocabulary, widen
or clear the retrieval scope (\S\ref{sec:scope}), or check the document
actually ingested (Library $\rightarrow$ does it appear, with the right
tags?).

\subsection{Conversations}

Every chat is saved automatically and appears under \textsc{Recent} in the
sidebar:

\begin{itemize}
  \item click a conversation to reopen it --- the full history, sources
        included;
  \item \emph{Search chats\ldots} filters the list as you type;
  \item the trash button next to \textsc{Recent} clears all conversations
        (confirmation asked);
  \item \emph{New chat} starts a fresh conversation --- the assistant carries
        recent turns of the \emph{current} conversation as context, so start
        a new chat when you change topic to keep answers sharp.
\end{itemize}

\subsection{Retrieval scope (searching within tags)}
\label{sec:scope}

Above the messages sits the \textbf{scope bar}. Expand it to see the tag
taxonomy grouped by type, and click tags to restrict retrieval: the assistant
will then search \textbf{only documents carrying at least one selected tag}
(OR semantics). The bar shows how many documents are currently in scope.

Typical uses: scope to \emph{Resumes} before asking ``which candidates know
Figma?'', or to a project tag before asking for a timeline summary. Clear the
scope to search the whole library again. Scope applies per conversation and
is sent with every question while active.

\subsection{The Document Library}

Sidebar $\rightarrow$ \textbf{Document Library}. This is the assistant's
knowledge base --- a table of every ingested document with its type, size,
tags, description, and upload info, plus a storage summary.

\begin{itemize}
  \item \textbf{Search documents\ldots} --- free-text filter on name and
        description.
  \item \textbf{Tag filter} --- click tag chips to filter the table (multiple
        tags narrow it down: a document must carry all of them).
  \item \textbf{Inline editing} --- each row lets you edit the document's
        tags and its description in place.
  \item \textbf{Delete} --- the row menu (\ldots) removes a document from the
        knowledge base (with confirmation). Its chunks leave the vector
        store; answers will no longer draw on it.
\end{itemize}

\subsection{Uploading documents}

Library $\rightarrow$ \textbf{Upload} opens the upload modal. Add individual
files or an entire \textbf{folder} (folder uploads group under the folder's
name and can share one set of tags and a description, applied to every file
inside).

\begin{itemize}
  \item \textbf{Supported formats:} PDF, DOCX, PPTX, XLSX, CSV, TXT, and
        Markdown.
  \item \textbf{Size limit:} 50\,MB per file. Oversized or unsupported files
        are flagged in the queue and skipped; remove them with the
        $\times$.
  \item \textbf{Tags at upload:} assign tags per file (or per folder) before
        uploading --- cheaper than retagging later.
  \item \textbf{Duplicates:} if the server already holds a file with the same
        name, the queue flags it and you resolve it inline:
        \textbf{Replace} the stored copy, \textbf{Keep both}, or
        \textbf{Remove} it from the queue.
\end{itemize}

Upload streams progress per file; each document is chunked, embedded, and
searchable as soon as its ingestion finishes.

\subsection{The tagging system}

Tags are the library's organising principle, and they do double duty: they
group and filter documents in the library, \textbf{and} they scope retrieval
in chat (\S\ref{sec:scope}).

The taxonomy has two levels, managed in Library $\rightarrow$ \textbf{Manage
Tags}:

\begin{itemize}
  \item \textbf{Tag types} --- categories such as \emph{Document kind},
        \emph{Project}, or \emph{Department}, each with a colour. A tag's
        colour comes from its type, so chips are visually scannable
        everywhere (library rows, scope bar, upload queue).
  \item \textbf{Tags} --- the values inside a type (e.g.\ \emph{Resume},
        \emph{Proposal}, \emph{Policy} under \emph{Document kind}).
\end{itemize}

In the Tag Manager you can create, rename, recolour, and delete types, and
add or delete tags within them. Deleting is confirmed first; a tag removed
from the taxonomy disappears from scope selection, and documents still
carrying an unknown tag show it in neutral grey.

\begin{gtip}
Design the taxonomy before bulk-uploading. A small, consistent set of types
(kind / project / team) beats dozens of ad-hoc tags --- scoped questions like
``summarise the \emph{Aurora} correspondence'' only work if the documents
were tagged consistently.
\end{gtip}

\section{Routine tasks and troubleshooting}

\subsection{Swapping the active model}

The stack ships several \texttt{qwen3} sizes trading speed for capability.
To switch:

\begin{itemize}
  \item \textbf{Panel:} Models view $\rightarrow$ pick the model
        $\rightarrow$ \emph{Use}.
  \item \textbf{CLI:} \texttt{./dffrnt-manager models use qwen3:8b}
\end{itemize}

Switching rewrites \texttt{llm\_model} in \texttt{config.toml} and restarts
the stack. To keep disk usage flat, the previously configured model is
removed from the cache --- except imported models, which are protected until
you switch to or delete them.

\textbf{On an air-gapped machine} the model must first be brought over as a
tarball:

\begin{enumerate}
  \item On any networked machine with this app:
\begin{lstlisting}
./dffrnt-manager models export qwen3:8b /media/usb/qwen3-8b.ollama.tar
\end{lstlisting}
        (pulls first if absent; the USB drive must not be FAT32 --- its
        4\,GB file limit is smaller than most models).
  \item On the air-gapped machine:
\begin{lstlisting}
./dffrnt-manager models import /media/usb/qwen3-8b.ollama.tar
./dffrnt-manager models use qwen3:8b
\end{lstlisting}
        (the import is checksum-verified).
\end{enumerate}

Both steps are also available in the panel's Models view.

\begin{gwarn}
The \emph{embedding} model is deliberately not swappable by command.
Changing it changes the vector dimension, which invalidates the entire index
and requires re-uploading or re-ingesting every document. If you must change
it, edit \texttt{embed\_model} \textbf{and} \texttt{vector\_size} in
\texttt{config.toml}, restart, and run \texttt{./dffrnt-manager reingest}.
\end{gwarn}

\subsection{Changing configuration}

Edit \texttt{config.toml} in the install folder (see
Annex~\ref{sec:annex} for every parameter), then:

\begin{lstlisting}
./dffrnt-manager restart
\end{lstlisting}

The config file survives reinstalls (only \texttt{install
--overwrite-config} replaces it, archiving the old file as
\texttt{config.toml.old}).

\subsection{Re-ingesting the library}

After changing chunking settings, or if a document seems poorly indexed:

\begin{lstlisting}
./dffrnt-manager reingest --dry-run   # preview what would be processed
./dffrnt-manager reingest             # asks for confirmation, then runs
\end{lstlisting}

Re-ingestion re-chunks and re-embeds every stored document from the original
uploaded files. It can take a while for large libraries (each chunk is
embedded by the local model).

\subsection{Backing up}

Everything lives in the install folder --- stop the stack, copy the folder,
start again:

\begin{lstlisting}
./dffrnt-manager stop
cp -r dffrnt /backup/dffrnt-$(date +%F)
./dffrnt-manager start
\end{lstlisting}

The critical subfolders are \texttt{data/} (original uploaded documents),
\texttt{qdrant\_storage/} (the vector index, tags, and saved conversations),
and \texttt{config.toml}. \texttt{ollama\_models/} is bulky and reproducible
(re-pull or re-import), so you may back it up less often.

\subsection{The audit trail}

\texttt{./dffrnt-manager audit} follows \texttt{logs/audit.jsonl}, the
application's append-only record of queries, uploads, and deletions ---
useful for reviewing how the knowledge base is being used on a shared
machine.

\subsection{Troubleshooting}

\paragraph{The first answer after a start is slow.}
On every start the models are pre-loaded into memory in the background; the
first question may land before warmup finishes. Subsequent answers run at
full speed. (On a fresh online install the very first warmup also waits for
the model download.)

\paragraph{Answers are slow in general.}
Model size vs.\ hardware. On CPU-only machines use a smaller model
(\texttt{models use qwen3:4b}). On NVIDIA machines make sure
\texttt{gpu = true} is set and the NVIDIA Container Toolkit is installed ---
the manager refuses GPU mode and says so when it isn't. On Apple Silicon the
GPU (Metal) is used automatically.

\paragraph{The assistant refuses although the answer is in a document.}
Check scope, rephrase with the document's vocabulary, confirm the document
ingested. If retrieval consistently misses, \texttt{top\_k} and
\texttt{score\_margin} in the annex are the tuning knobs (re-run the
evaluation suite before trusting changes).

\paragraph{Answers cite the wrong or outdated documents.}
Delete superseded documents from the library, or keep both but tag them so
scope can separate them. The system prompt prefers the latest agreed state
when documents contradict each other, but it can only work with what's
ingested.

\paragraph{An upload stays ``processing'' or fails.}
Check \texttt{./dffrnt-manager logs api}. Common causes: a corrupt or
password-protected file, or a file over 50\,MB. Fix the file and re-upload;
a failed file can simply be uploaded again.

\paragraph{A service is unhealthy / the UI won't load.}
\texttt{./dffrnt-manager status} to identify the service,
\texttt{logs <service>} to read its output, \texttt{restart} to recover. See
the Installer Guide for the full troubleshooting checklist (Docker down,
port conflicts, permissions).

\paragraph{Streaming looks ``chunky'' behind a reverse proxy.}
The API streams answers token by token and marks the responses no-buffering;
an intermediate nginx/CDN can still buffer them into bursts. Disable proxy
buffering for the \texttt{/api/*/stream} routes (details in the deployment
README).

%==============================================================================
\clearpage
\appendix
\renewcommand{\thesection}{\Alph{section}}

\section{Annex --- config.toml reference}
\label{sec:annex}

\texttt{config.toml} in the install folder is the single source of truth for
every tunable. Values can also be set as \texttt{UPPER\_SNAKE\_CASE}
environment variables; precedence is \textbf{environment $>$ config.toml $>$
built-in default}. After any change: \texttt{./dffrnt-manager restart}.

\begin{gwarn}
Parameters marked \warnmark\ invalidate the existing vector index ---
changing them requires re-uploading or re-ingesting every document.
\end{gwarn}

\subsection{Deployment}

\begin{pltable}{>{\ttfamily\footnotesize}p{3.6cm} >{\ttfamily\footnotesize}p{1.9cm} L}
\rowcolor{tablehead}\thd{Parameter} & \thd{Default} & \thd{Meaning}\\
\endfirsthead
\rowcolor{tablehead}\thd{Parameter} & \thd{Default} & \thd{Meaning}\\
\endhead
gpu & false & \normalfont Enable NVIDIA GPU acceleration for Ollama
(Linux/Windows container pathway). Requires the NVIDIA driver \textbf{and}
Container Toolkit on the host; the manager refuses to start GPU mode without
them. \textbf{Ignored on macOS}, where native Ollama uses Metal
automatically.\\
\end{pltable}

\subsection{Ollama --- models and generation}

\begin{pltable}{>{\ttfamily\footnotesize}p{3.6cm} >{\ttfamily\footnotesize}p{1.9cm} L}
\rowcolor{tablehead}\thd{Parameter} & \thd{Default} & \thd{Meaning}\\
\endfirsthead
\rowcolor{tablehead}\thd{Parameter} & \thd{Default} & \thd{Meaning}\\
\endhead
ollama\_url & \footnotesize http://localhost:11434 & \normalfont Host-side
URL of the Ollama server. Only the \textbf{port} is honoured across
deployment pathways --- keep the host \texttt{localhost}; the manager injects
the right in-container address itself.\\
llm\_model & qwen3:14b & \normalfont The chat LLM. Ships in three sizes:
\texttt{qwen3:4b} ($\sim$2.6\,GB, CPU-viable), \texttt{qwen3:14b}
($\sim$9.3\,GB, fits a 16\,GB GPU alongside the embedder),
\texttt{qwen3:30b-a3b} ($\sim$19\,GB, needs a 24\,GB GPU). Normally changed
via \texttt{models use}, not by hand.\\
embed\_model \warnmark & bge-m3 & \normalfont The embedding model that turns
text into vectors. Changing it changes the vector dimension --- deliberately
not switchable by command.\\
vector\_size \warnmark & 1024 & \normalfont Dimension of the embedding
vectors. \textbf{Must match the embed model's output dimension} (1024 for
bge-m3).\\
llm\_temperature & 0.1 & \normalfont Sampling temperature. Low = factual and
repeatable, which suits a grounded assistant; raise only if answers feel too
rigid.\\
llm\_timeout & 600 & \normalfont Seconds before a generation request is
abandoned. Local generation can legitimately be slow on big prompts.\\
llm\_num\_ctx & 16384 & \normalfont Context window in tokens. Ollama's own
default (4096) silently truncates RAG prompts once retrieved documents +
history grow; 16384 covers worst-case aggregate retrieval plus
\texttt{llm\_num\_predict}. Larger costs GPU memory ($\sim$1.4\,GB KV cache
at 16k for qwen3:14b).\\
llm\_top\_p & 0.95 & \normalfont Nucleus sampling cut-off (qwen3's
recommended value).\\
llm\_top\_k & 20 & \normalfont Sampling candidate cut-off (qwen3's
recommended value).\\
llm\_repeat\_penalty & 1.1 & \normalfont Penalises repetition. At
\texttt{1.0} the reasoning model can loop the same reasoning block forever on
ambiguous questions; \texttt{1.1} (qwen3's thinking-mode recommendation) lets
it converge.\\
llm\_num\_predict & 5000 & \normalfont Hard cap on generated tokens (thinking
+ answer) --- the runaway-generation safety net. Too low guillotines the
answer mid-thought; keep well under \texttt{llm\_num\_ctx} minus the prompt.
\texttt{-1} = unlimited (not recommended).\\
embed\_query\_prefix \warnmark & "" & \normalfont Task prefix prepended to
\emph{queries} before embedding (some models want
\texttt{"search\_query: "}). bge-m3 uses none --- leave blank.\\
embed\_document\_prefix \warnmark & "" & \normalfont Task prefix prepended to
\emph{documents} before embedding. bge-m3 uses none --- leave blank.\\
embed\_on\_cpu & false & \normalfont Pin the embedder to CPU. Set
\texttt{true} when the GPU can't hold the LLM and embedder at once --- keeps
the LLM resident instead of thrashing (evict-to-embed, reload-to-generate).
Ingestion gets slower; chat stays warm.\\
\end{pltable}

\subsection{Qdrant --- vector store and retrieval}

\begin{pltable}{>{\ttfamily\footnotesize}p{3.6cm} >{\ttfamily\footnotesize}p{1.9cm} L}
\rowcolor{tablehead}\thd{Parameter} & \thd{Default} & \thd{Meaning}\\
\endfirsthead
\rowcolor{tablehead}\thd{Parameter} & \thd{Default} & \thd{Meaning}\\
\endhead
qdrant\_url & \footnotesize http://localhost:6333 & \normalfont URL of the
Qdrant server.\\
collection\_name & \footnotesize dffrnt\_documents & \normalfont Collection
holding the document chunks.\\
summary\_collection\_name & \footnotesize dffrnt\_document\_ summaries &
\normalfont Collection holding per-document summaries (used for routing
aggregate queries).\\
tags\_collection\_name & \footnotesize dffrnt\_tags & \normalfont Collection
persisting the tag taxonomy.\\
conversations\_ collection\_name & \footnotesize dffrnt\_conversations &
\normalfont Collection persisting saved conversations.\\
distance & cosine & \normalfont Vector similarity metric: \texttt{cosine}
\textbar\ \texttt{dot} \textbar\ \texttt{euclid}. Leave at \texttt{cosine}
for bge-m3.\\
top\_k & 8 & \normalfont Retrieved-chunk candidates fetched per query.
Higher = wider evidence but longer prompts (mind
\texttt{llm\_num\_ctx}).\\
score\_margin & 0.20 & \normalfont Relative relevance cut: drop hits scoring
more than this below the best hit, so weak chunks never reach the prompt.
\texttt{0} disables. Corpus-dependent --- re-run the evaluation suite
(\texttt{pytest -m eval}) before changing.\\
aggregate\_group\_limit & 24 & \normalfont For aggregate (``all
candidates'', tag-scoped) queries: additionally fetch the best chunk(s) from
each of up to this many documents, merged into the candidate pool so
roll-ups cover every relevant document. \texttt{0} disables.\\
aggregate\_group\_size & 1 & \normalfont Chunks kept per document in that
grouped search.\\
\end{pltable}

\subsection{Chunking (\warnmark\ all four --- re-ingest after changing)}

\begin{pltable}{>{\ttfamily\footnotesize}p{3.6cm} >{\ttfamily\footnotesize}p{1.9cm} L}
\rowcolor{tablehead}\thd{Parameter} & \thd{Default} & \thd{Meaning}\\
\endfirsthead
\rowcolor{tablehead}\thd{Parameter} & \thd{Default} & \thd{Meaning}\\
\endhead
chunk\_strategy & recursive & \normalfont How documents are split:
\texttt{fixed} \textbar\ \texttt{sentence} \textbar\ \texttt{paragraph}
\textbar\ \texttt{recursive}. Recursive respects structure (paragraphs, then
sentences) while targeting the size below.\\
chunk\_size & 512 & \normalfont Target chunk size in characters. Small
chunks favour precise retrieval --- the priority for this application.\\
chunk\_overlap & 64 & \normalfont Characters shared between consecutive
chunks, so facts on a boundary aren't split away from their context.\\
chunk\_floor & 128 & \normalfont Chunks smaller than this are merged
forward, so headings and single-line fragments are never stored as useless
standalone chunks. \texttt{0} disables.\\
\end{pltable}

\subsection{Retrieval / prompting}

\begin{pltable}{>{\ttfamily\footnotesize}p{3.6cm} >{\ttfamily\footnotesize}p{1.9cm} L}
\rowcolor{tablehead}\thd{Parameter} & \thd{Default} & \thd{Meaning}\\
\endfirsthead
\rowcolor{tablehead}\thd{Parameter} & \thd{Default} & \thd{Meaning}\\
\endhead
history\_turns & 6 & \normalfont How many recent conversation turns
accompany each question as context. More = better follow-ups, longer
prompts.\\
system\_prompt & "" & \normalfont The assistant's standing instructions
(audience, tone, citation rules, refusal rule). Blank uses the built-in
default; the shipped \texttt{config.toml} carries DFFRNT's tuned prompt.
Edit with care --- the citation and refusal behaviour lives here.\\
\end{pltable}

\subsection{API / paths}

\begin{pltable}{>{\ttfamily\footnotesize}p{3.6cm} >{\ttfamily\footnotesize}p{1.9cm} L}
\rowcolor{tablehead}\thd{Parameter} & \thd{Default} & \thd{Meaning}\\
\endfirsthead
\rowcolor{tablehead}\thd{Parameter} & \thd{Default} & \thd{Meaning}\\
\endhead
api\_host & 0.0.0.0 & \normalfont Interface the API binds to.
\texttt{0.0.0.0} = every interface (needed inside the container).\\
api\_port & 8000 & \normalfont Host-side port for the API + UI. The
container always serves 8000 internally; this moves only the published
port.\\
data\_dir & data & \normalfont Where original uploaded files are stored
(relative to the install folder).\\
audit\_log\_path & \footnotesize logs/audit.jsonl & \normalfont Location of
the append-only audit trail.\\
\end{pltable}

\end{document}
